% s-216-p216-m004; original PDF 216, 217, 218
\Aphorism{s-216-p216-m004}{310}
Es giebt Dinge, die mit hundert Zungen reden und einen umfassenden Angriff, der kein Entrinnen lässt, auf das empfangende Subject führen; der Angegriffene hat gar nichts anderes zu thun, als sämmtliche Organe auszustrecken, sich mit sämmtlichen Sinnen an die Oberfläche zu begeben, damit ihm von der voll zuströmenden Mittheilung nichts verloren gehe. Derart reden manche berühmte Landschaften zur Seele des Beschauers, als Offenbarungen so polyphonen Stils und so brausender Instrumentation, dass Jeder wenigstens eine Stimme hören muss; ob das Ganze als Ganzes empfunden wird, ist sehr fraglich. In der Kunst ist aus dem Ideal vollständiger und erschöpfender Mittheilung das Princip des Gesammtkunstwerks abgeleitet worden; Orchester, Gesangsstimme, das Wort mit seinen Bild- und Gedankensuggestionen, der Mythos als eine Form des Primitiven, wobei dem reflexionsmüden modernen Menschen wohl wird, die Gebärde, das Costüm und endlich die Landschaft, die sympathische Resonanz des Seelenvorgangs — mit alledem wird der Empfangende zugleich bestürmt und derart in die Mitte genommen, dass er verstehen, miterleben \Sperr{muss}. Wagners Kunst ist strategisch; sie bricht Widerstände, sie behandelt den Hörer \Foreign{a priori} als Feind. — Es giebt noch eine andere Art Kunst und Natur, und bisweilen ist man besser auf sie als auf die anspruchsvollere, höhere Art gestimmt; ja, ich möchte sagen, so oft man gestimmt ist (und nicht erst gestimmt werden will), neigt man im Grunde einer bescheideneren Art Natur und Kunst zu: einer, die den Geniesser nicht vollständig zur receptiven Oberfläche ausspannt, ihn nicht ganz und gar auf den Empfänger-Posten befehligt, sondern ihm das Recht der Gegenrede einräumt. Ich denke an jene Erlebnisse, die nicht wie strömende Offenbarungen kommen, sondern langsam, fast widerwillig; an Landschaften ohne hervorstechenden Reiz, von einer schwer zu fühlenden Eigenart, die aber, einmal gefühlt, Eigenthum des Beschauers bleibt; an eine Kunst, die lieber zu wenig als zu viel thut, um verstanden zu werden, an deren seltsame, nicht im mindesten künstliche oder alterthümelnde Einfachheit man sich ungern, aber endlich doch gewöhnt und ihr zuletzt dafür dankt. Das Alpenthal, in dem ich dieses Lob der reizlosen, unaufdringlichen Kunst und Landschaft singe, ist selbst eine Landschaft solcher Art; es hat nicht die ewigen „malerischen Veduten“ der Berner Alpen; es schlägt den Betrachter nicht mit gigantischen Formen wie die Monte-Rosa-Gruppe: es schmeichelt sich zögernd ein mit schlichten, schlankgeformten Bergen und seinem ernsten Grün, und die Felszacken im Thalschluss, ohne den unvermeidlichen Gletscher, fallen nicht aus dem \Foreign{genus medium} in die grosse alpine Rhetorik, sondern geben nur leise den nächsten Contrast, die Dominante zur herrschenden Tonart an. — Wer sich mittheilt, will auch verstanden werden (die Symbolisten und die \Dots{3} Lyriker mit dem versagenden Athem ausgenommen), und eigentlich unzugängliche Dinge giebt es nicht im Bereiche der formen-, ton- und sprachbegabten Welt. Aber manches lässt sich mehr Zeit, verstanden zu werden, manches weniger, und wer zuviel giebt, kommt leicht in den Fall, gegen den Empfänger zudringlich zu erscheinen — „hat der Geber nicht zu danken, dass der Nehmende nahm?“
\par

